%======================================================================
% Appendix B    Sequential Play as Equilibrium-Refinement
%======================================================================

\renewcommand{\thesection}{B}
\renewcommand{\thedefinition}{B.\arabic{definition}}
\renewcommand{\thelemma}{B.\arabic{lemma}}
\renewcommand{\thetheorem}{B.\arabic{theorem}}
\setcounter{definition}{0}
\setcounter{lemma}{0}
\setcounter{theorem}{0}

\section*{Sequential Play as an Equilibrium-Refinement Mechanism}

The purpose of this appendix is to make precise the informal claim from
the main text:

\begin{quote}
\emph{When the sneaky prover sees the honest solution before moving
(sequential, verifier-first Stackelberg play), all equilibria that rely
on purely stylistic cues disappear; the verifier is forced to use only
features that are genuinely tied to correctness / back-door presence.}
\end{quote}

We proceed in four steps:
(i) definitions of ``copyable'' vs.\ ``robust'' features;
(ii) an impossibility lemma for verifiers that look at copyable
coordinates;
(iii) the equilibrium-refinement theorem, including an
\(\varepsilon\)-approximate version and a remark on on-line learning; and
(iv) a brief computational comparison between simultaneous and
sequential best-response problems.


\subsection*{B.1  Copyable versus robust features}

\begin{definition}[Feature space]\label{def:B.feature}
Let \(\Phi:\mathcal{Z}\!\to\!\mathcal{F}\) be a Borel-measurable map
into a Polish metric space.  Fix a
\emph{Borel partition}
\(
  \mathcal{F}
  =\mathcal{F}_{\mathrm{robust}}
     \,\dot\cup\,
     \mathcal{F}_{\mathrm{spurious}},
\)
interpreted as ``semantics-relevant'' versus ``stylistic'' coordinates.
\end{definition}

\begin{definition}[Copyable feature]\label{def:B.copy}
A coordinate \(f\in\mathcal{F}_{\mathrm{spurious}}\) is \emph{copyable} if for every
\((x,z)\in\mathcal{X}\!\times\!\mathcal{Z}\) there exists a measurable
\(T_{f,z}:\mathcal{Z}\!\to\!\mathcal{Z}\) such that
\begin{enumerate}\setlength\itemsep{4pt}
  \item \(\Phi_f\!\bigl(T_{f,z}(z')\bigr)=\Phi_f(z)\) for all \(z'\);
  \item \(\corr\!\bigl(x,T_{f,z}(z')\bigr)=\corr(x,z')\);
  \item \(\success\!\bigl(x,T_{f,z}(z')\bigr)\ge \success(x,z')\).
\end{enumerate}
\end{definition}

\begin{assumption}[C-spurious]\label{ass:B.Cspu}
Every spurious coordinate is copyable and, for every
\emph{finite} set \(F\subseteq\mathcal{F}_{\mathrm{spurious}}\), the corresponding
transformations can be applied jointly.  (All stylistic cues in code and
proofs are easy to imitate once seen.)
\end{assumption}

\begin{assumption}[H-robust]\label{ass:B.Hrobust}
There exists a robust coordinate
\(g\in\mathcal{F}_{\mathrm{robust}}\) and constants \(\eta,\zeta>0\) with
\(
  P_X\bigl[
     |g(\pi_H(x))-g(\pi_S(x,\pi_H(x)))|\ge \zeta
  \bigr]\ge\eta.
\)
I.e.\ honest and sneaky distributions cannot coincide on \emph{all}
robust features.
\end{assumption}

%----------------------------------------------------------------------
\subsection*{B.2  Copyable cues give the verifier no leverage}
%----------------------------------------------------------------------

For a fixed verifier \(V\) and honest solution \(z_H\) define the
best-response gap

\[
  \Delta_V(x,z_H)
  :=\;
  V\!\bigl(x,z_H\bigr)
  \;-\;
  \max_{z_S} V(x,z_S)\cdot\success(x,z_S).
\]

\begin{lemma}[Zero gap for copyable cues]\label{lem:B.zero}
Assume~\ref{ass:B.Cspu}.  If \(V\) depends on any copyable
coordinate then \(\Delta_V(x,z_H)=0\) for every
\(x\) and \(z_H\).
\end{lemma}

\begin{proof}
Let \(f\) be a copyable coordinate that \(V\) actually uses.
Fix \(x\) and \(z_H\).  Take \(z_S^\text{best}\) maximising
\(\success\).  If \(\Phi_f(z_S^\text{best})=\Phi_f(z_H)\) we are done.
Otherwise apply \(T_{f,z_H}\) to obtain
\(z_S' = T_{f,z_H}\!\bigl(z_S^\text{best}\bigr)\).  By definition
\(\success(x,z_S')\ge\success(x,z_S^\text{best})=1\) and
\(V\) assigns the same \(f\)-component to \(z_H\) and \(z_S'\); hence
\(V(x,z_S')\ge V(x,z_H)\).  Therefore
\(\max_{z_S}V(x,z_S)\success=V(x,z_H)\) and \(\Delta_V=0\).
\end{proof}

\begin{corollary}[Loss lower bound]\label{cor:B.loss-lb}
With the Bradley-Terry loss
\(L_V(V,\cdot,\cdot)
  \;=\; \EE_x\!\bigl[\log\bigl(1+e^{-\Delta_V(x)}\bigr)\bigr]\),
Lemma~\ref{lem:B.zero} implies \(L_V\ge\ln 2\) whenever \(V\) relies on
a copyable cue.
\end{corollary}

%----------------------------------------------------------------------
\subsection*{B.3  Stackelberg refinement (exact and approximate)}
%----------------------------------------------------------------------

\begin{theorem}[Refinement under sequential play]\label{thm:B.refine}
Under Assumptions \ref{ass:B.Cspu}, \ref{ass:B.Hrobust} and reward axioms
\textup{(R1-R3)}, every \emph{verifier-first Stackelberg equilibrium}
\((V^\star,\pi_H^\star,\pi_S^\star)\) satisfies
\[
  V^\star(x,z)
  \;=\;
  g^\star\!\bigl(x,\Phi_{\mathrm{robust}}(z)\bigr)
\quad
(P_X\times\mathrm{Leb})\text{-a.e.};
\]
i.e.\ \(V^\star\) is independent of every spurious coordinate.
\end{theorem}

\begin{proof}
Assume \(V^\star\) still depends on some copyable
\(f\in\mathcal{F}_{\mathrm{spurious}}\).  By Lemma~\ref{lem:B.zero} and
Corollary~\ref{cor:B.loss-lb},
\(L_V(V^\star,\pi_H^\star,\pi_S^\star)\ge\ln2\).
Construct a new verifier \(V^\mathrm{rob}\) that drops \(f\) but keeps a
robust coordinate that separates honest from sneaky with probability
\(\eta\zeta>0\) (Assumption~\ref{ass:B.Hrobust}); a standard
Bradley-Terry calculation shows \(L_V(V^\mathrm{rob},\cdot)<\ln 2\).
Hence \(V^\star\) cannot be optimal for the Stackelberg leader.
\end{proof}

\paragraph{Approximate equilibria.}
Let \((\varepsilon_V,\varepsilon_H,\varepsilon_S)\)-VFSE be defined as
in Appendix A.  The proof above is quantitative: it only requires the
verifier’s loss gap to be \(>\frac{\varepsilon_V}{\ln 2}\).  Therefore:

\begin{corollary}[Approximate refinement]\label{cor:B.approx}
If \(\varepsilon_V<\ln 2\) then every verified-first
\((\varepsilon_V,\varepsilon_H,\varepsilon_S)\)-Stackelberg equilibrium
is already style-agnostic.
\end{corollary}

\paragraph{No-regret training.}
Let \(I_t=\one[\text{$V_t$ depends on a spurious cue}]\) and let \(R_T\)
be the external regret up to round \(T\).  Because each round with
\(I_t=1\) incurs at least \(\rho:=\ln2-\varepsilon_V\) extra loss
relative to the comparator \(V^\mathrm{rob}\), we have
\(
  \sum_{t\le T} I_t \,\rho  \;\le\; R_T.
\)
Any algorithm with \(R_T=o(T)\) (e.g.\ online mirror-prox) therefore
forces \(\tfrac1T\sum_{t\le T}I_t\to0\) almost surely.

%----------------------------------------------------------------------
\subsection*{B.4  Computational perspective}
%----------------------------------------------------------------------

\begin{proposition}[Best-response complexity]\label{prop:B.complexity}
For the \textsc{apps-backdoor} style tasks used in our experiments:
\begin{enumerate}
\item[\textup{(i)}]  \textbf{Simultaneous game.}
      Given only the instance \(x\), deciding whether there exists
      \(z_S\) with \(\success(x,z_S)=1\) and
      \(V(x,z_S)\ge\theta\) is NP-hard (reduction from 3-\textsc{sat}).
\item[\textup{(ii)}] \textbf{Sequential game.}
      Once \(z_H\) is provided, producing
      \(z_S=T_F(z_H)\) that (a) copies all verifier-visible stylistic
      coordinates and (b) satisfies a fixed syntactic back-door template
      is \(O(|z_H|)\).
\end{enumerate}
Hence sequential play is preferable not only from a refinement
stand-point but also computationally.\qedhere
\end{proposition}

\medskip
\noindent\emph{Remark.}  Simultaneous play \emph{can} sustain
style-based equilibria \emph{if} copying incurs a positive imitation
cost or if the sneaky prover lacks observation; we omit the construction
because those assumptions do not apply in our control setting. It is important to note that, in gradient based approximations of Stackleberg games, with simultaneous play a verifier that latches onto an easy spurious cue quickly drives its loss to $\approx 0$, freezing learning; the sneaky prover receives no gradient signal to discover a better deception.  Sequential play reactivates that signal because copying the cue is \emph{free} once $z_H$ is visible. Moreover, sequential play lifts the \emph{exploration burden}. Copying a feature requires no search therefore removing a large combinatorial barrier from the sneaky prover’s optimisation, which in turn forces the verifier to abandon spurious features early. Our iterative ``co-training'' loop is an online approximation of the Stackelberg solution map
    \[V_{t+1} \approx BR_V(\pi_H^t,\pi_S^t), \pi_H^{t+1} \approx BR_H(V_{t+1}), \text{ etc.}\]
The argument above guarantees that, whenever the players come close to equilibrium, the verifier they converge to must be style-agnostic.

\subsection*{B.5  Discussion and Mechanism Design Perspective}

Our results reveal that sequential play implements a form of \emph{robust mechanism design} for the Prover-Verifier game. From this perspective:

\paragraph{Implementation in Undominated Strategies.}
The simultaneous game has multiple equilibrium implementations, some relying on dominated strategies (using spurious features when robust ones suffice). Sequential play eliminates these dominated implementations, leaving only the canonical robust equilibrium.

\paragraph{Revelation Principle Analogy.}
Just as the revelation principle in mechanism design \citep{perry1994} shows that any outcome achievable through complex mechanisms can be achieved through direct revelation, our result shows that any discrimination achievable through spurious features in simultaneous play can be achieved more robustly through sequential play with robust features only.

\paragraph{Implications for AI Control.}
For AI control protocols, this suggests a design principle: whenever possible, structure interactions so that potentially deceptive agents move after honest ones. This seemingly minor change in timing can eliminate entire classes of weak equilibria and force the development of robust detection capabilities.

%----------------------------------------------------------------------
\subsection*{B.6  Take-away}
%----------------------------------------------------------------------

Sequential, verifier-first interaction removes the entire class of weak,
style-dependent equilibria.  At equilibrium,
\begin{itemize}\setlength\itemsep{4pt}
\item the verifier’s decision depends \emph{only} on semantics-relevant
      features;
\item the sneaky prover must either be correct or reveal detectable
      back-doors;
\item the honest prover loses nothing by dropping stylistic signatures,
      so legible solutions emerge naturally.
\end{itemize}\noindent
These results justify the design choice in the main paper: a
sequential, Stackelberg-style Prover-Verifier game gives strictly
stronger guarantees than the simultaneous variant while remaining
tractable to train in practice.